\section{Desafios Sistêmicos para a Adoção no SUS}

A despeito do formidável potencial dos gêmeos digitais para a equidade na saúde bucal, a sua transição dos \textit{papers} acadêmicos para a realidade do Sistema Único de Saúde é obstaculizada por uma teia de desafios sistêmicos, operacionais e sociais. Não se trata meramente de superar limitações técnicas do software, mas de adequar uma infraestrutura de saúde continental a uma nova exigência computacional.

\subsection{Infraestrutura e ``Desertos Digitais''}
O obstáculo mais formidável é a assimetria na infraestrutura de conectividade. Estima-se que vastas porções das regiões Norte, Nordeste e áreas rurais do país permaneçam como autênticos ``desertos digitais''~\cite{cuadros2023assessing}. A transmissão de malhas 3D de alta densidade e de volumes tomográficos (DICOM) exige larguras de banda e estabilidade de rede que não estão disponíveis em milhares de UBSs. A adoção de tecnologias de computação de borda (\textit{edge computing}) e otimização de compressão, como o Structural Noise Scanner do OpenCode, torna-se crítica para operar sob condições de rede degradadas, mas não exime a necessidade de políticas estatais robustas de expansão do acesso à internet de alta velocidade.

\subsection{Custo de Aquisição de Hardware}
A construção de um gêmeo digital fidedigno depende da acurácia da captação inicial. Equipamentos como scanners intraorais de alta resolução (IOS) e tomógrafos de feixe cônico (CBCT) mantêm-se na faixa de dezenas a centenas de milhares de reais. Para a grande maioria dos municípios brasileiros — responsáveis pelo financiamento primário da atenção básica —, a aquisição e a manutenção continuada deste maquinário extrapolam a realidade orçamentária. A superação deste gargalo exige modelos inovadores de financiamento, como a criação de consórcios intermunicipais de saúde digital ou centros de diagnóstico móveis subsidiados pela esfera federal.

\subsection{Capacitação Profissional e Resistência Cultural}
A tecnologia só gera valor quando os operadores possuem proficiência técnica e confiança clínica no sistema. A transição do paradigma analógico (moldagens em alginato, radiografias 2D) para a aquisição digital exige programas massivos e contínuos de capacitação das equipes de saúde bucal. Além do desafio didático, observa-se uma barreira cultural: o receio da desumanização do atendimento e o ceticismo perante a ``caixa-preta'' algorítmica. O treinamento deve enfatizar que a inteligência artificial dos gêmeos digitais não substitui o julgamento clínico humano, mas atua como um sistema auxiliar de alta precisão (\textit{human-in-the-loop})~\cite{roy2025editorial}.

\subsection{Interoperabilidade e Fragmentação de Prontuários}
Por fim, o SUS sofre de uma crônica fragmentação informacional. Os sistemas de Prontuário Eletrônico do Paciente (PEP) frequentemente não dialogam entre as esferas municipal, estadual e federal, tampouco entre atenção primária e secundária. Um gêmeo digital odontológico fragmentado perde seu valor preditivo. É imperativa a adoção de padrões abertos de integração, como o FHIR (Fast Healthcare Interoperability Resources) associado a frameworks de gêmeos digitais modulares, como proposto pelo projeto OpenTwins~\cite{infante2025distributed}. Apenas um ecossistema semântico unificado garantirá que a história digital do paciente o acompanhe por toda a sua jornada terapêutica no sistema de saúde.
